%
% display-server.tex
%
% Z Specification for the Lux Display Server
% Formal model of DisplayServer._on_frame() and FrameReader
%

\documentclass[a4paper,10pt,fleqn]{article}
\usepackage[margin=1in]{geometry}
\usepackage{fuzz}
\usepackage[colorlinks=true,linkcolor=blue,citecolor=blue,urlcolor=blue]{hyperref}

\begin{document}

\title{Lux Display Server: A Z Specification}
\author{Formal Model of Frame Lifecycle, Client Management, and Scene State}
\date{March 2026}
\hypersetup{
  pdftitle={Lux Display Server: A Z Specification},
  pdfauthor={Punt Labs},
  pdfsubject={Z Specification},
  pdfcreator={fuzz/probcli}
}
\maketitle

\tableofcontents
\newpage

% ===================================================================
\section{Introduction}
% ===================================================================

The Lux display server is an ImGui render loop that listens on a Unix
domain socket and renders scenes.  Each frame, it accepts new
connections, polls clients for messages, renders the current scene,
and flushes interaction events.

This specification models the stateful core:
\begin{itemize}
  \item Client connection tracking and FrameReader association
  \item Scene replacement, incremental patching, and clearing
  \item Message dispatch (scene, update, clear, ping)
  \item Interaction event generation and broadcast
  \item The FrameReader streaming buffer
\end{itemize}

We abstract away ImGui rendering (it is stateless from the
protocol's perspective) and focus on the state that the protocol
cares about.

% ===================================================================
\section{Basic Types}
% ===================================================================

\begin{zed}
  [FD, ELEMID, SCENEID, RAWBYTES, ACTION]
\end{zed}

$FD$ models socket file descriptors.  $ELEMID$ and $SCENEID$ are
opaque identifiers for elements and scenes.  $RAWBYTES$ represents
a chunk of bytes received from a socket.  $ACTION$ identifies
the action string on a button click.

% ===================================================================
\section{Free Types}
% ===================================================================

% Boolean type (fuzz doesn't have built-in boolean)
\begin{zed}
  ZBOOL ::= ztrue | zfalse
\end{zed}

Element kinds supported in Phase~1:
\begin{zed}
  ElementKind ::= ekText | ekButton | ekSeparator | ekImage
\end{zed}

Client message types:
\begin{zed}
  ClientMsgType ::= cmScene | cmUpdate | cmClear | cmPing
\end{zed}

Display message types (sent back to clients):
\begin{zed}
  DisplayMsgType ::= dmReady | dmAck | dmPong | dmInteraction
\end{zed}

% ===================================================================
\section{Global Constants}
% ===================================================================

\begin{axdef}
  maxClients : \nat \\
  maxElements : \nat \\
  maxEvents : \nat \\
  maxBufSize : \nat
  \where
  maxClients = 64 \\
  maxElements = 1000 \\
  maxEvents = 256 \\
  maxBufSize = 16777216
\end{axdef}

% ===================================================================
\section{Element}
% ===================================================================

An element within a scene.  We model only the identity and kind;
the rendering attributes (content, label, path, etc.)\ are opaque
to the protocol state machine.

\begin{schema}{Element}
  elemId : ELEMID \\
  elemKind : ElementKind
\end{schema}

% ===================================================================
\section{Scene}
% ===================================================================

A scene is an ordered list of elements with a unique scene ID.
Element IDs within a scene must be unique.

\begin{schema}{Scene}
  sceneId : SCENEID \\
  elements : \seq Element
  \where
  \# elements \leq maxElements \\
  \forall i, j : \dom elements @
\quad~ i \neq j \implies
\quad~ (elements~i).elemId \neq (elements~j).elemId
\end{schema}

% ===================================================================
\section{FrameReader}
% ===================================================================

The FrameReader accumulates raw bytes from a socket and tracks
how many complete messages are available for draining.

We abstract the buffer as a natural number (bytes accumulated)
and track how many complete messages have been parsed.

\begin{schema}{FrameReaderState}
  bufSize : \nat \\
  pendingMsgs : \nat
  \where
  bufSize \leq maxBufSize
\end{schema}

\begin{schema}{InitFrameReaderState}
  FrameReaderState'
  \where
  bufSize' = 0 \\
  pendingMsgs' = 0
\end{schema}

Feeding bytes into the reader increases the buffer size.

\begin{schema}{FeedBytes}
  \Delta FrameReaderState \\
  bytesIn? : \nat
  \where
  bytesIn? > 0 \\
  bufSize + bytesIn? \leq maxBufSize \\
  bufSize' = bufSize + bytesIn? \\
  pendingMsgs' = pendingMsgs
\end{schema}

Draining extracts all complete messages, reducing the buffer.
The number of messages drained is non-deterministic (depends on
message boundaries within the byte stream).

\begin{schema}{DrainMessages}
  \Delta FrameReaderState \\
  drained! : \nat \\
  bytesConsumed? : \nat
  \where
  bytesConsumed? \leq bufSize \\
  bufSize' = bufSize - bytesConsumed? \\
  pendingMsgs' = 0 \\
  drained! = pendingMsgs + bytesConsumed?
\end{schema}

% ===================================================================
\section{Display Server State}
% ===================================================================

The server tracks connected clients, their FrameReaders, the
current scene (if any), and a queue of pending interaction events.

\begin{schema}{DisplayServer}
  clients : \power FD \\
  readers : FD \pfun FrameReaderState \\
  hasScene : ZBOOL \\
  currentScene : Scene \\
  eventQueue : \seq ELEMID \\
  listening : ZBOOL
  \where
  \dom readers = clients \\
  \# clients \leq maxClients \\
  \# eventQueue \leq maxEvents \\
  hasScene = ztrue \implies
\quad~ (\forall eid : \ran eventQueue @
\quad~\quad~ eid \in \{ e : \ran (currentScene.elements) @ e.elemId \})
\end{schema}

The key invariants are:
\begin{enumerate}
  \item Every client has exactly one FrameReader ($\dom readers = clients$).
  \item Client count is bounded.
  \item Queued events reference elements in the current scene.
\end{enumerate}

% ===================================================================
\section{Initialization}
% ===================================================================

\begin{schema}{InitDisplayServer}
  DisplayServer'
  \where
  clients' = \emptyset \\
  readers' = \emptyset \\
  hasScene' = zfalse \\
  eventQueue' = \langle \rangle \\
  listening' = ztrue
\end{schema}

\begin{schema}{Init}
  InitDisplayServer
\end{schema}

% ===================================================================
\section{Operations}
% ===================================================================

% -------------------------------------------------------------------
\subsection{Accept Connection}
% -------------------------------------------------------------------

A new client connects.  We allocate a fresh FrameReader for it.

\begin{schema}{AcceptConnection}
  \Delta DisplayServer \\
  newClient? : FD
  \where
  listening = ztrue \\
  newClient? \notin clients \\
  \# clients < maxClients \\
  clients' = clients \cup \{ newClient? \} \\
  (\exists frs : FrameReaderState |
\quad~ frs.bufSize = 0 \land frs.pendingMsgs = 0 @
\quad~ readers' = readers \cup \{ newClient? \mapsto frs \}) \\
  hasScene' = hasScene \\
  currentScene' = currentScene \\
  eventQueue' = eventQueue \\
  listening' = listening
\end{schema}

% -------------------------------------------------------------------
\subsection{Disconnect Client}
% -------------------------------------------------------------------

A client disconnects (error or clean close).

\begin{schema}{DisconnectClient}
  \Delta DisplayServer \\
  deadClient? : FD
  \where
  deadClient? \in clients \\
  clients' = clients \setminus \{ deadClient? \} \\
  readers' = \{ deadClient? \} \ndres readers \\
  hasScene' = hasScene \\
  currentScene' = currentScene \\
  eventQueue' = eventQueue \\
  listening' = listening
\end{schema}

% -------------------------------------------------------------------
\subsection{Receive Scene}
% -------------------------------------------------------------------

A client sends a new scene, replacing any existing one.
The event queue is cleared since old element references are stale.

\begin{schema}{ReceiveScene}
  \Delta DisplayServer \\
  newScene? : Scene
  \where
  hasScene' = ztrue \\
  currentScene' = newScene? \\
  eventQueue' = \langle \rangle \\
  clients' = clients \\
  readers' = readers \\
  listening' = listening
\end{schema}

% -------------------------------------------------------------------
\subsection{Clear Scene}
% -------------------------------------------------------------------

A client sends a clear message, removing all content.

\begin{schema}{ClearScene}
  \Delta DisplayServer
  \where
  hasScene' = zfalse \\
  eventQueue' = \langle \rangle \\
  clients' = clients \\
  readers' = readers \\
  listening' = listening
\end{schema}

% -------------------------------------------------------------------
\subsection{Remove Element (Patch)}
% -------------------------------------------------------------------

An update message removes an element by ID from the current scene.

\begin{schema}{RemoveElement}
  \Delta DisplayServer \\
  targetId? : ELEMID
  \where
  hasScene = ztrue \\
  targetId? \in \{ e : \ran (currentScene.elements) @ e.elemId \} \\
  currentScene'.sceneId = currentScene.sceneId \\
  \ran (currentScene'.elements) =
\quad~ \{ e : \ran (currentScene.elements) | e.elemId \neq targetId? \} \\
  \# currentScene'.elements = \# currentScene.elements - 1 \\
  hasScene' = hasScene \\
  eventQueue' = eventQueue \\
  clients' = clients \\
  readers' = readers \\
  listening' = listening
\end{schema}

% -------------------------------------------------------------------
\subsection{Button Click (Interaction)}
% -------------------------------------------------------------------

The display detects a button click.  An interaction event is queued
for broadcast to all clients.

\begin{schema}{ButtonClick}
  \Delta DisplayServer \\
  buttonId? : ELEMID
  \where
  hasScene = ztrue \\
  buttonId? \in \{ e : \ran (currentScene.elements) |
\quad~ e.elemKind = ekButton @ e.elemId \} \\
  eventQueue' = eventQueue \cat \langle buttonId? \rangle \\
  \# eventQueue < maxEvents \\
  hasScene' = hasScene \\
  currentScene' = currentScene \\
  clients' = clients \\
  readers' = readers \\
  listening' = listening
\end{schema}

% -------------------------------------------------------------------
\subsection{Flush Events}
% -------------------------------------------------------------------

At the end of each frame, queued events are broadcast to all
clients and the queue is emptied.

\begin{schema}{FlushEvents}
  \Delta DisplayServer
  \where
  eventQueue' = \langle \rangle \\
  hasScene' = hasScene \\
  currentScene' = currentScene \\
  clients' = clients \\
  readers' = readers \\
  listening' = listening
\end{schema}

% -------------------------------------------------------------------
\subsection{Shutdown}
% -------------------------------------------------------------------

The display server shuts down: all clients are disconnected,
the scene is cleared, events are flushed.

\begin{schema}{Shutdown}
  \Delta DisplayServer
  \where
  clients' = \emptyset \\
  readers' = \emptyset \\
  hasScene' = zfalse \\
  eventQueue' = \langle \rangle \\
  listening' = zfalse
\end{schema}

% ===================================================================
\section{System Invariants}
% ===================================================================

The key properties maintained across all operations:

\begin{enumerate}
  \item \textbf{Reader--client bijection}: Every connected client
    has exactly one FrameReader, and every FrameReader belongs to
    a connected client ($\dom readers = clients$).

  \item \textbf{Event validity}: Queued interaction events reference
    elements that exist in the current scene.

  \item \textbf{Bounded resources}: Client count, element count,
    event queue length, and buffer sizes are all bounded.

  \item \textbf{Scene atomicity}: A new scene completely replaces
    the old one and clears stale events.

  \item \textbf{Unique element IDs}: Within a scene, no two elements
    share an ID.
\end{enumerate}

\end{document}
